(1) 单位质量能量 $E=(\dot x^2+\omega(t)^2x^2)/2$。沿解求导 $\dot E=\dot x(\ddot x+\omega^2x)+\omega\dot\omega x^2=\omega\dot\omega x^2$。频率调制本身可作功，因而瞬时机械能一般不守恒；这一精确式不依赖绝热近似。

(2) 设 $x=A(t)\cos\phi(t)$、$\dot\phi=\omega(t)$。将表达式代入方程，在 $|\dot\omega|\ll\omega^2$ 下保留一级慢变量，令快正弦项系数为零，得到 $2\omega\dot A+\dot\omega A=0$，因此 $A(t)=A_0\sqrt{\omega(0)/\omega(t)}$，$\phi(t)=\phi_0+\int_0^t\omega(s)ds$。

(3) 一个局域振动周期内 $\langle x^2\rangle\simeq A^2/2$、$\langle\dot x^2\rangle\simeq\omega^2 A^2/2$。于是 $\langle E\rangle\simeq A^2\omega^2/2$，再使用 $A^2\omega\simeq\mathrm{const}$，得到 $J=\langle E\rangle/\omega\simeq A_0^2\omega(0)/2$。也能将精确 $\dot E$ 取周期平均复核该结论。

(4) 当 $\dot\omega=0$，上述表达式退化为精确谐振解 $x=A_0\cos(\omega_0t+\phi_0)$，$J$ 严格守恒。若 $\omega$ 接近零、频率变化与周期同量级、或快变参数触发共振，$|\dot\omega|/\omega^2\ll1$ 失效，公式不再保证近似守恒。
